\icode{E-EME-SI-L3-261121-003}
在如图所示的空间直角坐标系中，一不计重力且带正电的粒子从坐标为(\(L\),0,\(L\))处以某一初速度平行\(y\)轴方向射出，经时间\(t\)，粒子前进的距离为\(L\)。在该空间加上匀强电场，粒子仍从同一位置以相同的速度射出，经相同时间\(t\)后恰好运动到坐标原点\(O\)。已知粒子的比荷为\(k\)，则该匀强电场的场强大小为\paren[]

\begin{tikzpicture}[>=stealth,font=\scriptsize]
	\draw[->,blue] (0,0)--(0,2)node[right]{\(z\)};
	\draw[->,blue] (0,0)--(2,0)node[above]{\(y\)};
	\draw[->,blue] (0,0)--(-1.2,-1.2)node[right]{\(x\)};
	\draw[dashed] (-.8,-.8)node[right]{\(L\)}--(-.8,.4)--(0,1.2)node[right]{\(L\)};
	\draw[->,thick] (-.8,.4)--(-.2,.4);
	\fill (-.8,.4) circle(.05);
	\node[blue] at(.1,-.2){\(O\)};
\end{tikzpicture}
\begin{choices}
	\item \(\dfrac{2\sqrt{3}L}{kt^2}\)
	\item \(\dfrac{2\sqrt{5}L}{kt^2}\)
	\item \(\dfrac{2L}{kt^2}\)
	\item \(\dfrac{4L}{kt^2}\)
\end{choices}
\begin{answer}
	
\end{answer}